\subsection{World creation and computational-value curvature}
\label{sec:world-value}
A generative architecture does not select a welfare functional, but a free
margin in the number of worlds changes the value of additional computation.
To expose that margin, let $C_{\rm tot}\geq u>0$ be service available to
worlds, let each world receive $c\geq u$, and initially relax world count
to a real variable $W\geq1$. The constraint is $Wc=C_{\rm tot}$. The
exponent $\theta\in[0,1]$ values world population; the elasticity
$\alpha\in(0,1]$ values resources within a world. This distinction is
between two margins of valuation, not between simulated and physical
persons.

\begin{lemma}[A free world-number margin]
\label{lem:world-curvature}
For the continuous relaxation,
\begin{equation}
 \max_{Wc=C_{\rm tot},\,W\geq1,\,c\geq u}W^\theta c^\alpha
 =\begin{cases}
 u^{\alpha-\theta}C_{\rm tot}^{\theta},&\theta>\alpha,\\
 C_{\rm tot}^{\alpha},&\theta\leq\alpha.
 \end{cases}
 \label{eq:world-value}
\end{equation}
Thus the value exponent is $v=\max(\theta,\alpha)>0$. Writing
$C_{\rm tot}=u(1+C)$ and normalizing the value at the mandatory root
as $[V(u(1+C))/V(u)-1]/v$ gives exactly $F_{1-v}(C)$.
For the separate specification $W^\theta\ln(c/u)$, $\theta=0$ gives
$\ln(C_{\rm tot}/u)$, while $\theta>0$ gives
$(C_{\rm tot}/u)^\theta/(\exp(1)\theta)$ only once
$C_{\rm tot}\geq u\exp(1/\theta)$.
\end{lemma}
\noindent\emph{Proof sketch.}
Substitution of $W=C_{\rm tot}/c$ reduces the power case to the sign
of $\alpha-\theta$. In the logarithmic case, optimize
$\exp[\theta\ln(C_{\rm tot}/u)-\theta v_c]v_c$ over
$0\leq v_c\leq\ln(C_{\rm tot}/u)$. The full proof, ties, and the
integer-world correction are in \cref{app:world-proof}.

At $\theta=\alpha$ every feasible continuous split ties. With integer
world counts and $\theta>\alpha$, the exact value becomes
$C_{\rm tot}^\alpha\lfloor C_{\rm tot}/u\rfloor^{\theta-\alpha}$;
it has the same asymptotic exponent, not the same differentiable function.
A regular recursive hierarchy restricts the available counts further but
preserves this exponent. The shifted smooth family is therefore exact for
the stated relaxation, not for discrete launches or indivisible worlds.

Average-log valuation is one interpretation of the baseline $\eta=1$.
It is not forced by computational sophistication. Total-log valuation
with freely increasing world count instead has linear asymptotic value,
and the power model with $\theta=1$ has linear value exactly. At
$\eta=0$ the corner condition becomes
\begin{equation}
 p_ja\ell'(B)\geq r_j+S_j.
 \label{eq:totalist-corner}
\end{equation}
For logarithmic local production, densification must satisfy
$a/(1+B)\geq\max_j(r_j+S_j)/p_j$. More valued lives can thus strengthen
the attraction of unused external resources rather than eliminate it.
This connects the opportunity-cost comparison to
\citet{Bostrom2003Waste}, without attributing the present production
function or equation to that paper.

Recursion supplies one implementation of an open-ended world-number
margin. Its analytical role is now explicit: admitting that margin and
specifying its population weight changes the induced value curvature.
A non-recursive implementation with the same feasible world counts has
the same implication. Recursion neither manufactures additional substrate
nor chooses $\theta$ on behalf of a civilization. The extensive world
margin, rather than nesting by itself, is the operative assumption.
